\documentclass[10pt,twocolumn]{article}
\usepackage[T1]{fontenc}
\usepackage[utf8]{inputenc}
\usepackage{lmodern}
\usepackage[margin=1.8cm,top=2.4cm,bottom=2cm,columnsep=0.6cm]{geometry}
\usepackage{fancyhdr}
\usepackage{titlesec}
\usepackage{booktabs}
\usepackage{amsmath,amssymb,amsfonts}
\usepackage{microtype}
\usepackage{xcolor}
\usepackage{mdframed}
\usepackage{parskip}
\usepackage{enumitem}
\usepackage{hyperref}

\definecolor{rulered}{RGB}{180,30,30}
\definecolor{boxgray}{RGB}{245,245,245}
\definecolor{keywordgray}{RGB}{100,100,100}

\pagestyle{fancy}
\fancyhf{}
\fancyhead[L]{\textsc{\small Claude Mag}}
\fancyhead[C]{\small\textit{Machine Learning Research Digest}}
\fancyhead[R]{\small 2026-08-12}
\fancyfoot[C]{\small\thepage}
\renewcommand{\headrulewidth}{0.8pt}

\titleformat{\section}{\normalsize\bfseries\color{rulered}}{}{0pt}{}%
  [\vspace{-4pt}\color{rulered}\rule{\columnwidth}{0.4pt}]
\titlespacing{\section}{0pt}{8pt}{4pt}

\hypersetup{colorlinks=true,urlcolor=rulered,linkcolor=black}

\begin{document}
\setlength{\parindent}{0pt}
\setlength{\parskip}{4pt}

{\small\color{keywordgray}\textbf{Keywords:}
  domain adaptation \textbullet{}
  time series \textbullet{}
  path signatures \textbullet{}
  distribution alignment}\par\vspace{4pt}

{\LARGE\bfseries\raggedright
  CPDA: Class-Conditional Path Distribution
  Alignment for Unsupervised Time-Series
  Domain Adaptation}\par\vspace{4pt}

{\large\itshape\raggedright
  A Composite Signature-Spectral Kernel Beats 30 Baselines
  on 13 Time-Series DA Benchmarks Without Adversarial Training}\par\vspace{6pt}

{\small\textbf{By} Felix Ott, Christopher Mutschler}\par\vspace{2pt}

{\small\color{keywordgray}arXiv:2608.09193 \textbullet{} Preprint, August 2026}
\par\vspace{2pt}\noindent{\color{rulered}\rule{\columnwidth}{0.6pt}}\vspace{6pt}

Most unsupervised time-series domain adaptation (UTS-DA) methods align only global
marginal feature distributions, discarding the class structure that defines decision
boundaries. CPDA instead aligns class-conditional latent \emph{path distributions}
via a novel composite signature-spectral kernel that jointly captures temporal
trajectory structure, frequency content, and semantic features. Evaluated on
13 benchmarks against 30 baselines, CPDA achieves consistent state-of-the-art
results with CNN, ResNet18, and TCN backbones---without adversarial training.

\begin{mdframed}[backgroundcolor=boxgray,linecolor=rulered,linewidth=1pt,
    innerleftmargin=6pt,innerrightmargin=6pt,
    innertopmargin=6pt,innerbottommargin=6pt]
\textbf{\small AT A GLANCE}\par\vspace{4pt}
\begin{description}[leftmargin=0pt,labelsep=4pt,itemsep=2pt,parsep=0pt,
    font=\small\bfseries]
  \item[Problem:] UTS-DA distribution shifts arise from different users, sensors,
    or temporal dynamics; prior methods align only global marginals, losing class
    structure.
  \item[Why it matters:] Labeling data for every new deployment context is
    infeasible; UTS-DA enables transfer without target labels in IoT, health,
    and manufacturing.
  \item[Method:] Class-conditional alignment via a composite kernel combining
    path signatures, spectral features, semantic pooling, and low-rank dynamics.
  \item[Result:] State-of-the-art on 13 benchmarks against 30 baselines with
    3 backbone architectures.
  \item[Evidence:] Non-adversarial; stable training across all experimental
    configurations.
\end{description}
\end{mdframed}

\section{The Big Picture}

Time-series classifiers deployed in the real world face a pervasive problem:
the distribution of training data differs from the distribution at deployment.
A human activity recognition model trained on data from a 25-year-old male with
one wearable device may perform poorly on a 60-year-old female with a different
device. A machine fault-detection system trained on one production line may fail
on a nominally identical line with slightly different resonance characteristics.

Unsupervised time-series domain adaptation addresses this by transferring knowledge
from a labeled \textbf{source domain} to an unlabeled \textbf{target domain}.
The key insight of CPDA is that temporal data has rich structure---\textbf{path
trajectories} in feature space, \textbf{frequency spectra}, and \textbf{class-specific
distributional signatures}---that global marginal alignment discards. CPDA preserves
this structure through a principled kernel design.

\section{Why Existing Approaches Fall Short}

The dominant paradigm in UTS-DA follows the deep domain adaptation tradition from
computer vision: minimize a distributional discrepancy (e.g., MMD) or train a domain
adversary (e.g., DANN) on marginal feature distributions. Both approaches have
structural limitations for time-series data.

\textbf{Global marginal alignment} ignores class structure. If source class~A
aligns with target class~B due to superficial distributional similarity, a classifier
that looked correct in source space will fail in target space. This \textbf{negative
transfer} effect is well-known but remains underaddressed in UTS-DA.

\textbf{Spectral and temporal structure} is discarded when features are pooled
before alignment. A model that misses the difference between a 10~Hz and a 15~Hz
vibration pattern may align the two classes' pooled statistics while failing on
the actual classification problem.

\textbf{Adversarial training} is unstable, sensitive to hyperparameters, and prone
to mode collapse, making it difficult to deploy reliably across diverse benchmark
settings.

\section{Core Mechanism}

CPDA operates in the encoder's latent space, treating each sample as a
\textbf{path}---a sequence of latent feature vectors over time. It aligns
class-conditional path distributions via a composite kernel $k_{\mathrm{CPDA}}$
that combines four components:

\begin{enumerate}[itemsep=2pt,parsep=0pt]
  \item \textbf{Semantic pooling kernel:} Classic RBF kernel on time-averaged
    latent features; captures overall distributional similarity.
  \item \textbf{Path-signature kernel:} Measures similarity of ordered polynomial
    statistics of the time-series trajectory---captures temporal dynamics.
  \item \textbf{Spectral kernel:} Fourier-based kernel on the frequency-domain
    representation of the latent path---captures periodic structure.
  \item \textbf{Low-rank path-signature kernel:} Compressed, dominant modes of
    temporal variation via truncated signature expansion---reduces noise sensitivity.
\end{enumerate}

\section{Technical Formulation}

Let $\phi_\theta: \mathcal{X} \to \mathbb{R}^{T \times d}$ be the encoder mapping
a time-series $x$ to a $T$-step latent path. Define the class-conditional source
and target distributions:
\[
P^y_S = \{\phi_\theta(x) \mid x \in \mathcal{D}_S, y_x = y\}, \quad
P^y_T = \{\phi_\theta(x) \mid x \in \mathcal{D}_T, \hat{y}_x = y\}
\]
where $\hat{y}_x$ are pseudo-labels for target samples.

The CPDA discrepancy is:
\[
\mathcal{D}_{\mathrm{CPDA}} = \sum_{y=1}^{C} \mathrm{MMD}^2_{k_{\mathrm{CPDA}}}(P^y_S, P^y_T)
\]
with:
\[
k_{\mathrm{CPDA}}(\gamma, \gamma') =
  \alpha_1 k_{\mathrm{sem}}(\gamma, \gamma') +
  \alpha_2 k_{\mathrm{sig}}(\gamma, \gamma') +
  \alpha_3 k_{\mathrm{spec}}(\gamma, \gamma') +
  \alpha_4 k_{\mathrm{lr}}(\gamma, \gamma')
\]
where $\gamma, \gamma'$ are latent paths and $\alpha_i$ are learned mixture weights.

The training objective is:
\[
\mathcal{L} = \mathcal{L}_{\mathrm{CE}}(\mathcal{D}_S) +
              \lambda\,\mathcal{D}_{\mathrm{CPDA}}(\mathcal{D}_S, \mathcal{D}_T)
\]

\section{Learning or Inference Procedure}

\textbf{Pseudo-label generation:} Target pseudo-labels are produced by the
source-trained classifier and updated iteratively as the encoder improves.
High-confidence pseudo-labels (above a threshold) are used for alignment.

\textbf{Optimization:} Stochastic gradient descent on $\mathcal{L}$; no adversarial
loop, no discriminator. The MMD-based objective is smooth and differentiable.

\textbf{Inference:} Standard forward pass through the adapted encoder and classifier;
no test-time computation overhead relative to the base model.

\section{What the Guarantee Says}

Under standard domain adaptation assumptions (covariate shift, shared label
conditional), minimizing class-conditional MMD drives the target encoder toward
the source class-conditional distributions, providing a theoretical basis for
classifier transfer. The path-signature kernel is characteristic on path space
under mild regularity conditions, guaranteeing that $\mathcal{D}_{\mathrm{CPDA}} = 0$
implies equality of path distributions. In practice, finite-sample approximation
introduces error, but the rich kernel reduces the approximation gap relative
to simpler kernels.

\section{Experimental Findings}

\begin{itemize}[itemsep=2pt,parsep=0pt]
  \item \textbf{13 benchmarks:} Including activity recognition (HAR, UCIHAR,
    WISDM), fault detection (CWRU, PU), and physiological monitoring (EEG, ECG).
  \item \textbf{30 baselines:} Spanning discrepancy-based (MMD, CORAL),
    adversarial (DANN, CDAN, AdvSKM), and pseudo-labeling methods (SHOT, SDAT).
  \item \textbf{3 backbones:} CNN, ResNet18, TCN---CPDA achieves state-of-the-art
    with all three.
  \item \textbf{Margin:} CPDA outperforms the second-best method on average
    across all 13 benchmarks by a consistent margin with all backbone choices.
\end{itemize}

\section{Ablations and Interpretation}

\textbf{Kernel component ablation.}
Removing any single component ($k_{\mathrm{sig}}$, $k_{\mathrm{spec}}$,
$k_{\mathrm{lr}}$, $k_{\mathrm{sem}}$) degrades performance on most benchmarks,
confirming that each component captures complementary distributional information.
The path-signature kernel has the largest marginal contribution on benchmarks
with strong temporal dynamics (fault detection); the spectral kernel dominates
on benchmarks with periodic signals (EEG, ECG).

\textbf{Global vs.\ class-conditional alignment.}
Replacing class-conditional alignment with global marginal MMD using
$k_{\mathrm{CPDA}}$ degrades performance substantially on benchmarks with
closely adjacent class clusters, confirming the negative-transfer hypothesis.

\textbf{Non-adversarial vs.\ adversarial.}
CPDA trains stably across all 13 benchmarks; GAN-based methods (DANN, CDAN)
show higher variance and failed convergence on 3 of 13 benchmarks.

\section*{Reference}

Felix Ott, Christopher Mutschler.
\textbf{CPDA: Class-Conditional Path Distribution Alignment for
Unsupervised Time-Series Domain Adaptation}.
arXiv:2608.09193, August 2026.
\url{https://arxiv.org/abs/2608.09193}

\end{document}
